% =====================================================================
%  FICHE 11 — THEME 3 — Corrigé FACILE — Équilibrer une équation redox
% =====================================================================
\documentclass[11pt,a4paper]{article}
\usepackage[utf8]{inputenc}
\usepackage[T1]{fontenc}
\usepackage{lmodern}
\usepackage[french]{babel}
\usepackage{amsmath,amssymb}
\usepackage[version=4]{mhchem}
\usepackage{siunitx}
\sisetup{output-decimal-marker={,},per-mode=symbol}
\usepackage{xcolor}
\usepackage{array}
\usepackage{booktabs}
\usepackage{enumitem}
\usepackage[most]{tcolorbox}
\usepackage[a4paper,margin=2.1cm,headheight=26pt,headsep=10pt]{geometry}
\usepackage{fancyhdr}
\usepackage[hidelinks]{hyperref}

\definecolor{primary}{RGB}{146,95,160}
\definecolor{primarydark}{RGB}{92,52,104}
\definecolor{corrcol}{RGB}{0,135,90}
\newcommand{\NO}{\ensuremath{\mathrm{N.O.}}}

\newcounter{exo}
\newtcolorbox{corr}[1][]{enhanced,breakable,boxrule=0.9pt,arc=2pt,colback=corrcol!4,
  colframe=corrcol,fonttitle=\bfseries,coltitle=white,left=3mm,right=3mm,top=2.2mm,bottom=2.2mm,
  title={Corrigé~\refstepcounter{exo}\theexo\ifx\relax#1\relax\relax\else\ \textmd{--- \itshape #1}\fi}}

\pagestyle{fancy}\fancyhf{}
\fancyhead[L]{\small\textcolor{primarydark}{\textbf{Fiche 11 · Thème 3} — Corrigé}}
\fancyhead[R]{\small\textcolor{primarydark}{\textsc{Facile}}}
\fancyfoot[C]{\small\thepage}
\renewcommand{\headrulewidth}{0.8pt}
\renewcommand{\headrule}{\hbox to\headwidth{\color{primary}\leaders\hrule height \headrulewidth\hfill}}

\begin{document}
\begin{center}
{\LARGE\bfseries\color{primarydark}Thème 3 — Corrigé Facile}\\[-2pt]
{\color{corrcol}\rule{\textwidth}{1pt}}
\end{center}

\begin{corr}[la demi-équation du permanganate]
Mn : $+\mathrm{VII}\to+\mathrm{II}$, donc $n=5$.
\begin{itemize}[leftmargin=1.5em,topsep=1pt,itemsep=1pt]
  \item \textbf{Oxygène} : 4 O à gauche (\ce{MnO4^-}), 0 à droite $\to$ ajouter 4 \ce{H2O} à droite.
  \item \textbf{Hydrogène} : $4\times2=8$ H à droite $\to$ ajouter 8 \ce{H+} à gauche.
  \item \textbf{Charge} : à gauche $-1+8=+7$, à droite $+2$ $\to$ ajouter 5 \ce{e^-} à gauche.
\end{itemize}
\[
\boxed{\ce{MnO4^- + 8 H+ + 5 e^- -> Mn^{2+} + 4 H2O}}
\]
\end{corr}

\begin{corr}[le dichromate — attention au facteur 2]
Cr : $+\mathrm{VI}\to+\mathrm{III}$ ($\Delta=3$ par Cr), mais \ce{Cr2O7^{2-}} contient \textbf{2} Cr.
\begin{itemize}[leftmargin=1.5em,topsep=1pt,itemsep=1pt]
  \item \textbf{Chrome} : 2 Cr à gauche $\to$ \textbf{2} \ce{Cr^{3+}} à droite. Électrons : $2\times3=\mathbf{6}$.
  \item \textbf{Oxygène} : 7 O à gauche $\to$ 7 \ce{H2O} à droite.
  \item \textbf{Hydrogène} : $7\times2=14$ H $\to$ 14 \ce{H+} à gauche.
  \item \textbf{Charge} : gauche $-2+14=+12$, droite $2\times(+3)=+6$ $\to$ 6 \ce{e^-} à gauche. \checkmark
\end{itemize}
\[
\boxed{\ce{Cr2O7^{2-} + 14 H+ + 6 e^- -> 2 Cr^{3+} + 7 H2O}}
\]
\end{corr}

\begin{corr}[le nitrate en monoxyde d'azote]
N : $+\mathrm{V}\to+\mathrm{II}$, $n=3$.
\begin{itemize}[leftmargin=1.5em,topsep=1pt,itemsep=1pt]
  \item Oxygène : 3 O à gauche, 1 à droite $\to$ 2 \ce{H2O} à droite.
  \item Hydrogène : 4 H à droite $\to$ 4 \ce{H+} à gauche.
  \item Charge : gauche $-1+4=+3$, droite $0$ $\to$ 3 \ce{e^-} à gauche.
\end{itemize}
\[
\boxed{\ce{NO3^- + 4 H+ + 3 e^- -> NO + 2 H2O}}
\]
\end{corr}

\begin{corr}[deux demi-équations sans oxygène]
\[
\ce{Fe^{2+} -> Fe^{3+} + e^-} \quad(\text{Fe : }+\mathrm{II}\to+\mathrm{III},\ \textbf{oxydation}) ; \qquad
\ce{Sn^{2+} -> Sn^{4+} + 2 e^-} \quad(\text{Sn : }+\mathrm{II}\to+\mathrm{IV},\ \textbf{oxydation}).
\]
Comme le \NO\ augmente et que les électrons sont \emph{produits} (à droite), ce sont bien deux oxydations.
\end{corr}

\begin{corr}[le contrôle des charges]
$\ce{MnO4^- + 8 H+ + 5 e^- -> Mn^{2+} + 4 H2O}$ :
\begin{itemize}[leftmargin=1.5em,topsep=1pt,itemsep=1pt]
  \item Charges à gauche : $(-1)+8(+1)+5(-1)=-1+8-5=+2$. À droite : $+2 + 0 = +2$. \checkmark
  \item Mn : 1 = 1 ; O : 4 = 4 ; H : $8 = 4\times2 = 8$. \checkmark
\end{itemize}
La demi-équation est donc \textbf{correctement équilibrée} (atomes \emph{et} charges).
\end{corr}

\end{document}
